\chapter{隔离记录}
\chapterkana{隔離}
\epigraphbox{“形如槁木，心如死灰。”}{——《庄子·齐物论》}

\kaishou{前一九九〇年那一击没有要她的命。\par{}五年以后，他们把她接了回去，关进一间不跟六环城相通的隔离舱，一关就是一千七百六十四年。\par{}账本上一年只写两个字，「如故」；写到最后一年，才换了四个字。}

\section{接回}

纯狐启动那台遗留设备之后，档案在公元前一九九〇年前后只留下一行注有「未解」的记录；紧接着的五年里，月面没有留下任何下文。直到公元前一九八五年，另外一批档册上才重新出现地之主任的行踪：她被从地面接回月面，安置进一处与六环城彼此隔绝的封闭舱段。\footnote{接回之事亦见传世《淮南子·览冥训》“羿请不死之药于西王母，姮娥窃以奔月”，刘安撰，刘文典《淮南鸿烈集解》，中华书局 1989 年点校本，第 211 页。传世文献记作私奔，档案则记作专程接回。}

五年之间到底发生过什么，目前没有任何文字材料可以作证。也许月之都在等待地面设备彻底停机，也许接回行动本身需要重新开凿不经过六环城内部的专用通道；无论是哪一种情况，五年这段空白本身就说明，接回是一次经过掂量之后的善后，跟当场抢救无关。读者若把前一章未了的那一条线头拉到此处，便会看见那次善后在档案里留下了一套独立于《章程总目》之外的记录册。

这一章的主证物是《隔离舱管理档》，编入基金会档案时编号 \texttt{ФВ-1969-03}，一共存下来十四块薄合金板。合金板长约二十二厘米、宽十四厘米、厚不足一毫米，材质成分与第三章所见的《章程总目》属同一路冶炼工艺，只是版面裁得更窄，左侧打有两个直径三毫米的穿线孔，原本应当用细金属丝串成一叠。十四块板子正反两面都密密麻麻划着细格，每一格只有黄豆大小，刻刀笔锋比《总目》还要细上一半。第十四块的右下角已经蚀断，断口呈斜向崩裂，最后几栏只剩半行字迹，其余各板保存尚算完整。\footnote{十四块板的编目附注说明，同箱另有专用通道的维修零片三件，尺寸残缺过甚，难以识读，此处不引。}

从舱段的构造和进出规程来看，这套隔离设施沿用了第五章提过的净秽《出界》旧例，却又多加了一道锁。舱门不朝六环城内开启，所有空气、给水与日用补给全走单向管道，守门与送档的人在舱外隔着观察窗登记，期满并不放回居住环。从形式上看，这更像一座专为一个人长期预备的封闭舱，与城内那套“居外舍三月乃入环”的临时隔离截然分属两回事。

\section{一千七百六十四年}

《隔离舱管理档》的体例极简，效法《岁计》卷的成例，按城内年序一年占一格，平时只在「状貌」与「言动」两栏里填两个字：「如故」。有变化的年份只在极少数格子里留下一行短批。

\begin{archive}[内部参考]{档案 № ФВ-1969-03}
第一板首格（折合前一九八五年）：「创面不合，色如常。」\par{}
第一板第四格（折合前一九八八年）：「问其名，不答。」\par{}
第二板首格（折合前一九六〇年前后）：「围增三寸。」\par{}
第十四板末格（折合前二二一年）：「状貌：酷似蟾蜍之物。」其下版缘断缺。
\end{archive}

要把这些城内年序折算回地面公元纪年，靠的是前后两头对校：起点接在第八章考定的公元前一九九〇年纯化变故之后，终点则是前二二一年撤离前最后一次清点。\footnote{城内年序与地面公元纪年的换算限制，见《月之都——〈广寒宫残简〉辑佚与考释》附录二；本书所用前 1985 年是据前 221 年末格上溯一千七百六十四栏推算所得。}一头一尾算下来，这份档案整整记了一千七百六十四年。绝大部分格子里千篇一律写着「如故」，真正有字句变动的只有四处，其中三处最能说明问题。

第一处是接回当年的首格：「创面不合，色如常。」不合，是指那处伤口没有重新长拢；色如常，是指边缘既不溃烂，也不变黑坏死。两半批注摆在一处，恰恰把第八章末提出的疑团证了一半：被纯狐那台设备拆解过的那一小片组织，既不在自身修复的名单里，也没有沦为会腐败的死肉，它只是被彻底划出了这具身体的管理范围之外。

第二处落在前一九六〇年前后，批作「围增三寸」。此后每隔几十年至百余年，这一栏里就会重现一次「围增」，数目大小不等，增生的方向一律朝外。把第二板到第十四板的尺度数字从头到尾加在一起，增长极慢：一千年出头一共只多出了不到四尺，然而整整一千多年里从未真正停过。

第三处也是最触目的一处，留在第十四板断口之前的末格：「酷似蟾蜍之物。」这一栏平常只省事地写「如故」，到前二二一年撤离前夕突然改了措辞，写完这六个字之后，板面便告中断。

把这三处连在一起，看见的便是一道缓慢移动的边界。被纯化掉的那一片组织已经不再属于她自己，机体的修复程序找不到可作用的对象，只好贴着边界停住；然而界线外侧的躯体仍保有再生能力，每次试图往创口修补时，新长出来的组织照样以那片死界为参照，随即被残留的界线一并划进去。界线不仅沿着体表往四周扩散，也朝内部推进，走过的地方一圈叠着一圈堆在创面边缘。一千年里增出来的那不到四尺，就是这样一层一层堆出来的；堆到前二二一年，原初的轮廓已经完全走样，记档的人才会写下「酷似蟾蜍之物」那六个字。

\section{不老与不死}

写到这里，第八章留到本章才作答的那个问题就躲不过去了：如果她已经失去了一块不再属于自己的身体，而且创面还在千年之内缓慢外推，为什么她没有像一般人那样死在地面，反而在隔离舱里一直活到了前二二一年。回答这个问题，必须先把「不老」和「不死」拆开来看。前几章通过《章程总目》和《名籍》卷证明过，六环城里的居民并不老病，第四百零九枚背面那人连续当差三百二十七年也没有告老；但同一份《名籍》卷里写得很清楚，城中仍有「失足」除籍的先例，第十章将考证的乱事里城中一次便死了两万多人。这意味着靠净秽系统隔离了变化的月之民，做到的是不老、不病，却从来没有做到不死；一旦遭遇外伤或高强度的毁坏，照样会立刻毙命，而且受伤之后不会自己长好。

真正能在一千七百六十四年的外伤与边界蚕食下既不老又不死的，在整套世界观与档案里只指向第六章出现过的那东西：蓬莱药。第六章已经看过，净秽山早期囤积的是把城内污秽抽压之后无害化制成的「蓬莱药」。当时档案与钻孔只说明那是一种停工后才出事的库存，到了这一章，才能看出这种产物若被活人直接服下会发生什么后果。

凡是服过蓬莱药的人，在档案里称作「服药者」，后世多叫作蓬莱人。药使人既不衰老也不死亡：身体受创之后会重新长合，哪怕碎裂也能慢慢复原。然而这份不死带有一笔极高的账要付。第五章说过，所谓污秽在月之都的体例里就是变化本身；服药者维持着永不熄灭的再生与生命过程，等于成了一个关不掉的变化源头，从此昼夜不停地向周围产出污秽。

这一层一旦落定，第四章与第六章留下的两条旧账便同时合上了口。两笔账一笔关乎当年为何下地，一笔关乎今日为何不死，拆开看各站半边，合起来正好拼实。

第一，嫦娥早在公元前二〇五〇年下地之前，就已经是服过蓬莱药的服药者。正因为她持续向外产出污秽，才无法继续留在讲究无秽的六环城内，地面于她与其说是外派的驻地，不如说是一座天然的流放处；月之民把她送下去的行动，在档案措辞里写作「用」，而不用寻常外派官员的「派」字。

第二，纯狐的纯化之所以能在一名服药者身上留下不可逆的创口，并没有推翻蓬莱药本身的药效。设备解除的是被处理那一小片躯体与宿主之间的归属，使那一片脱离蓬莱人之身，不再受再生机制管辖；其余尚未被划入界线的部分依旧照常服从药效，照样不老不死。两样加在一起，才酿出隔离舱里那一幕怪状：整个人死不了，那一片却再也长不回原样，连新长出的组织也跟着一并堆成了畸形。

\section{忘却与不答}

除了「状貌」一栏之外，第十四块金属板与第一块金属板之间，还有一条一直写到中断都没有改口的话，出现在「言动」栏。与年年写「如故」的常规记录不同，这一栏记下的是隔离者对外部问询的反应。

公元前一九八八年，也就是她被接回隔离舱的第三年，第一板第四格记下了一行短语：「问其名，不答。」此后几十年间，隔三岔五还能看到「再问，不答」的字样；再往后，连「问」的记录都逐渐稀了下去，只剩年复一年的「如故」。

这四个字放在月之都的制度背景下看，格外值得留意。第六章曾经提过，月之民停止自身变化之后，原本在地面与近地环境里起作用的记忆消除场就不再对他们生效。\footnote{关于无秽状态、变化停止与记忆消除场之间的关系，参见索科洛夫《记忆消除场通论》。本书只取档案所见的行为后果，不探讨场本身的物理作用机制。}城里的人因此成了绝少遗忘的一群：从登船起经历过的事，两万年来件件都在脑子里摆得整整齐齐，只是轻重缓急的尺度早已冻死在离开地面的那一天，记住一切，却不再觉得有任何一件事重要。

隔离舱里的嫦娥正好站在同一条规则的反面。她服过蓬莱药，体内一刻不停地进行着变化，也在一刻不停地产生污秽，无法像城里人那样冻结记忆权重；那个对六环城内十余万人全部失灵的记忆消除场，在她身上于是照常生效，并不因污秽多少而增减，只是把她拉回与常人无异的受场状态，而在无限延长的寿限里一轮一轮冲刷下来，遗忘便无休无止地累积。\footnote{持续产生污秽只是解除权重冻结，并未放大场强；漫长遗忘出自常效之场在无限寿限下的反复冲刷，详见索科洛夫《记忆消除场通论》。}她在隔离舱里不断忘掉刚发生不久的事。城还在的时候，尚有按时前来查验舱段的吏员，把旧卷、名字、往事一遍遍摊在她眼前，她还能顺着眼前这些纸片暂时捡回一部分记忆；等到前二二一年整座城断电撤离，门外再也没有人来翻旧账，她便只能一直忘下去：忘掉历法是怎么通过自己送进中原的，忘掉自己为何被锁在舱里，忘掉自己叫什么名字，到头来连维持生命最起码的举动也会渐渐失落，偏偏又死不了。

于是同一条规则在隔着一道舱壁的两边，裂成了两种截然相反的结局：一边是记得住两万年前每一笔账、却什么都不再在乎的一城人，一边是把一切忘得干干净净、却永远死不掉的一个人。两种截然对立的境遇都由同一套无秽与变化之辨生出来，中间并没有另设两套规矩。

不过，「问其名，不答」未必全由忘却所致。站在一个被当作流放品送往地面、经手发出过灭门处置令、随后又被锁回独立舱段的人的立场上，面对舱外那些隔窗登记的吏员保持沉默，同样说得通。前文把这一句归入常态遗忘重新起效的早发迹象，已经掺进了一半推论；至于当年隔窗开口问她名字的那名吏员，后来在撤离前是否还记得她的本名，十四块板子上同样没有交代。

\section{五个母题}

如果把《隔离舱管理档》和传世文献里关于嫦娥奔月的叙事摆在一处，传世故事里反复出现的五个核心母题，几乎都能在这十四块合金板上找到对应的落点。为了避免空口比附，现将两头的条目照直列在表 \ref{tab:chang-e-motifs} 里。\footnote{五项神话母题主要见《淮南子·览冥训》高诱注及《初学记》《太平御览》所引《灵宪》残文，成书跨汉晋数百年，细节互有增饰。}

\begin{table}[htbp]
  \centering
  \caption{嫦娥故事五个母题与隔离档案对照}
  \label{tab:chang-e-motifs}
  \begin{tabular}{@{}p{18mm}p{40mm}p{66mm}@{}}
    \toprule
    传世母题 & 常见叙事 & 本书据档案的读法 \\
    \midrule
    奔月 & 服药飞升、入月宫 & 前 1985 年自地面接回月面 \\
    窃药 & 盗取西王母所赠不死药 & 混淆「西王母」所赠历法与自身早先服药 \\
    蟾蜍 & 入月后化为蟾蜍 & 第十四板末格「酷似蟾蜍之物」 \\
    捣药 & 月中兔终年捣药为之治病 & 月兔本职为加工污秽，与隔离者无关 \\
    囚禁 & 孤独幽居广寒宫中 & 独立封闭舱段，与六环城不连通 \\
    \bottomrule
  \end{tabular}
  \tabsource{传世母题据《淮南子·览冥训》及《后汉书·天文志》注引张衡《灵宪》；档案读法据 \texttt{ФВ-1969-03} 及第六章所引净秽山材料。}
\end{table}

这五个对应里，最容易被后世误会的是「捣药」与「窃药」两项。第四章和第六章已经拆解过，月兔捣制蓬莱药是城中维持生态循环的例行杂役，远在后羿代夏、嫦娥接回之前就已经天天在做，只是后世听见天上既有孤居的嫦娥、又有捣药的白兔，顺理成章便把两件不相干的事拧成了一股。至于「窃药」一节，更像是两件年代相差五十年的旧事在口传里叠在了一处：前二〇五〇年下地之前她就已经在月面服下了药，前二〇〇〇年那位被地面记作「西王母」的月面传历者下地送来的，只是一部删改过的历法（地面误传为西王母赠药），后世既把历法误作仙丹，又把她最终被接回月面误作吞丹升天。至于这位在当地留下「西王母」名号的传历者在月之都档册里究竟叫什么名字，目前尚无直接档案可呼应，要留到第十二章合参日本列岛早期出土的木简与口传材料后才能正式对号。

\section{断角之外的空白}

这一章讲完了，照例要回头把撑不住的地方拆开来看。手头的实物只有那十四块薄合金板，外加传世文献里几段彼此抵牾的旧故事；连隔离舱本身在月面的具体坐标、供电如何接往地核链路，都只靠《章程总目》的只言片语和第六章的钻探报告往回补齐。

退一步核查全章。先要承认的是，没有去过月面的人对着这十四块隔海运来的合金板做推断，原本就比现场勘查隔着几层隔阂。

具体到正文各步，有三处悬在半空。第一处就是前一九八八年那句「问其名，不答」：本书顺着遗忘重启的规则，把它看作记忆流失重新开始的证据，可老实说，这一步是猜的；如果她当年仅仅是不愿开口理会关押自己的人，档案上的四个字写出来也一模一样。第二处是公元前一九九〇年到一九八五年之间那五年空档：档案既没记清谁去地面接人，也没说清用什么载具带回一个正在持续释放污秽的活体。第三处也是最紧要的一处，藏在十四块板子那上千个「如故」里。前二二一年末格虽然写下了「酷似蟾蜍之物」，但中间横跨一千六百多年的大半格子里连一次尺码重测都没有留下；在封闭舱段长期无事可记的情况下，值守吏员到底是一年一年真开窗核查后才写「如故」，还是接连几十年懒得掀盖、直接照抄上一格的旧字，直到撤离前才草草看了一眼并补上最后六个字，单从金属板自身的刻痕深浅已经分不出来了。若中间那一千多年里少有人掀开隔离窗，本书依据逐年累加推出来的「不到四尺」线性速率，便要重新打个折扣。

第十四块合金板在公元前二二一年的「酷似蟾蜍之物」下面陡然断了口子，隔离舱自身的记录到此戛然而止。不过，为什么这座城偏偏会在这一年全体撤空、又为什么早在公元前一九八〇年之后山里就不再有新药入库、只剩原状污秽一直堆到最后，单看一个隔离舱是答不上来的。前五章和第六章挂在半空中的「月兔之乱」与「处置就地」，还有第七章里只露了个名字的「天文密葬法」，都要留待最后三章依次考辨并收拢。\footnote{关于传历者「西王母」在月面档案中的本名及其东渡始末，将在第十二章合参日本列岛早年木简与口传材料后正式考定。}
